\section{随机变量的数字特征}
在应用中，我们往往感兴趣于某些能描述随机变量特征的常数。这种常数被称为随机变量的\textbf{数字特征}。对于随机变量，一些经典的数值特征是期望（均值），方差，标准差，协方差，矩。

\subsection{均值}

对任意随机变量 $X: \Omega \to \mathbb{R}$，定义其\textbf{均值}为：
\[
\E[X] = \int_\Omega X(\omega) \, \dif P(\omega).
\]

这是个从测度角度的定义，我们也可以给出简单一点的连续型随机变量$X$的均值定义为：

\[
\E[X]=\int_{-\infty}^{\infty} xf_X(x)\dif x
\]

离散型就是把积分改为求和，就是换了个测度。注意到我们有处理无限积分/和的情况，所以有时均值不一定存在。

它的意义基本就是一个按概率的加权平均数。

对于一个存在均值的随机变量$X$，我们把$X-\E[X]$叫做\textbf{中心化}的随机变量

均值是对于单个随机变量而言的，不过对于随机变量函数，我们也可以计算其均值。为了方便讨论，我们只研究连续随机变量函数。

\begin{theorem}
设 \(\mathbf{X} = (X_1, X_2, \dots, X_n)^T\) 是 \(n\) 维连续型随机向量，其联合概率密度函数为 \(f_{\mathbf{X}}(\mathbf{x})\)。设 \(g: \mathbb{R}^n \to \mathbb{R}\) 是一个可测函数，定义新的随机变量 \(Y = g(\mathbf{X})\)。
\[
\E[g(\mathbf{X})] = \int_{\mathbb{R}^n} g(\mathbf{x}) f_{\mathbf{X}}(\mathbf{x}) \, \dif \mathbf{x}
\]
\end{theorem}

\begin{proof}
由定义
\[
\E[g(\mathbf{X})] = \int_\Omega g(\mathbf{X}(\omega))  \dif P(\omega)
\]

由于 $g$ 可测，应用积分变换：
\[
\int_\Omega g(\mathbf{X}(\omega)) \dif P(\omega) = \int_{\mathbb{R}^n} g(\mathbf{x}) \dif P_{\mathbf{X}}(\mathbf{x})
\]

再应用概率密度的定义：
\[
\int_{\mathbb{R}^n} g(\mathbf{x}) \dif P_{\mathbf{X}}(\mathbf{x}) = \int_{\mathbb{R}^n} g(\mathbf{x}) f_{\mathbf{X}}(\mathbf{x}) \dif\mathbf{x}
\]

综上即得
\[
\E[g(\mathbf{X})] = \int_{\mathbb{R}^n} g(\mathbf{x}) f_{\mathbf{X}}(\mathbf{x}) \, \dif\mathbf{x}
\]
\end{proof}

一个直观的理解就是我们可以直接把这个函数拿到积分内。对于线性函数来说，应用定理得到
\begin{align*}
E\left[\sum_{i=1}^n a_i X_i + b\right]
&= \int_{\mathbb{R}^n} \left( \sum_{i=1}^n a_i x_i + b \right) f_{\mathbf{X}}(\mathbf{x}) \dif\mathbf{x} \\
&= \left(\sum_{i=1}^n\int_{\mathbb{R}^n}   a_i x_i  f_{\mathbf{X}}(\mathbf{x}) \dif\mathbf{x} \right)
   + \int_{\mathbb{R}^n} b f_{\mathbf{X}}(\mathbf{x}) \dif\mathbf{x}\\
&= \sum_{i=1}^n a_i E\left[X_i\right] +b 
\end{align*}
这就是期望的线性性。它还告诉我们：常数的均值是常数。注意线性性这一事实并不依赖于随机变量的任何独立性。

随机变量的独立性则体现在它们乘积的期望上。

\begin{align*}
\E[XY] &= \int_{\mathbb{R}^2} xy f_{XY}(x,y) \dif x\dif y\\
\E[X]\E[Y] &= \left(\int_{\mathbb{R}} x f_{X}(x) \dif x\right) \left(\int_{\mathbb{R}} y f_{Y}(y)\dif y\right)\\
&= \int_{\mathbb{R}^2} xy f_{X}(x)f_{Y}(y) \dif x\dif y
\end{align*}

可见当$X$与$Y$独立时，上下两个式子相等。

\begin{attention}
上面的结论反过来是不一定成立的。$\E[X]\E[Y]=\E[XY]$只能得出$f_{X}(x)f_{Y}(y)=f_{XY}(x,y)$\textbf{几乎处处}相等。
\end{attention}

\subsection{方差}

随机变量$X$的\textbf{方差}定义为$\E[(X-\E[X])^2]$。

连续型随机变量$X$的\textbf{方差}定义为：

\[
\Dev[X]=\int_{-\infty}^{\infty} (x-\E[X])f_X(x)\dif x
\]

离散型就是把积分改为求和。注意到我们有处理无限积分/和的情况，所以有时方差也可能不存在。

它描述了随机变量的分布偏离均值的情况，并且总是非负的。

由于对于常数$C$，$\E[C]=C$，我们可以得到$\Dev[C]=0$。这一结论反过来也\textbf{几乎}是成立的，但我们把证明留在弱大数定律的部分，我们先放心使用这个结论。请放心，弱大数定律的证明不会依赖这中间的任何东西。

对于一般的随机变量函数，它的方差没有什么太好的结论，只有在应用线性函数时才有比较好的结论。不过为了叙述的简洁性，我们把它的性质留到协方差这一部分来介绍。

\textbf{标准差}就是方差的平方根。

\subsection{协方差}

对于两个随机变量$X$，$Y$，它们的协方差定义为

\[
\Cov(X,Y) = \E\left[(X-\E[X])(Y-\E[Y])\right] = \E[XY]-\E[X]\E[Y]
\]

按式子就有：$\Dev[X]=\Cov(X,X)=\E[X^2]-(\E[X])^2$。我们还可以验证这一表达式的对称性和双线性性：
\[
\Cov(X,Y)=\Cov(Y,X)
\]
\[
\Cov(\alpha X_1+\beta X_2,Y) = \alpha \Cov(X_1,Y) +\beta \Cov(X_2,Y)
\]

同时还有平移不变性$\Cov(X+C,Y)=\Cov(X,Y)$，和非负性$\Cov(X,X)\geq 0$。

这些性质代入到方差中，就可以知道：

\[
\Dev[\lambda X] = \lambda^2\Dev[X],\quad \Dev[X+C]=\Dev[X]
\]

所有这些性质的证明基本都依赖于均值的线性性质。

在适当的函数空间上，我们可以把协方差看作一个内积。如果我们把两个相差一个\textbf{几乎处处}为常值随机变量的随机变量视为等价的（相当于我们在处理中心化后的随机变量），那么协方差就有\textbf{正定性}了，也就构成一个实内积了。标准差于是就是它诱导的范数。

相应地，我们可以给出内积经典结论。

首先是勾股定理
\begin{align*}
\Dev[X+Y]&=\Cov(X+Y,X+Y)\\
&= \Cov(X,X) + \Cov(Y,Y) + 2\Cov(X,Y)\\
&= \Dev[X]+\Dev[Y]+2\Cov(X,Y)
\end{align*}

于是$\Dev[X+Y]=\Dev[X]+\Dev[Y]$当且仅当$\Cov(X,Y)=0$。

这种“协方差正交性”我们称为\textbf{不相关性}。

然后是柯西不等式：
\[
\Cov(X,Y)^2\leq \Dev[X]\Dev[Y]
\]

等号成立当且仅当一个是另一个的标量倍。注意这里我们是在零均值意义下讨论的，所以它们还可能相差一个常数。不过总归是$Y=a+bX$这种线性关系就是了。

这里很自然地，这里就有一个类似余弦的量：

\[
\rho_{XY} = \frac{\Cov(X,Y)}{\sqrt{\Dev[X]\Dev[Y]}} \leq 1
\]

这就称为\textbf{相关系数}。反过来协方差也可以表示成$\rho_{XY}\sqrt{\Dev[X]\Dev[Y]}$。

由前面的式子可以知道，$X,Y$独立时有$\Cov(X,Y)=\E[XY]-\E[X]\E[Y]=0$。也就是说：\textbf{独立性蕴含不相关性}。但不相关的变量不一定独立。我们只能通过换位得到并非不相关的变量不是独立的。

对于一组随机变量$\mathbf{X}=(X_1,\ldots,X_n)^T$，我们也可以给出一个类似Gram矩阵的构造：协方差矩阵。

我们称矩阵$(C)_{ij} = \Cov(X_i,X_j)$为它们的协方差矩阵，也记作$\Cov(\mathbf{X})$。这是个对称半正定矩阵，我们可以轻松用谱定理分解它。如果我们丢掉小特征值对应的部分，就能对数据进行降维分析。如果对$\mathbf{X}$做一个线性变换得到$\mathbf{Y}=A\mathbf{X}+\mathbf{b}$，那么有$\Cov(\mathbf{Y})=A\Cov(\mathbf{X})A^T$

\subsection{矩}
矩是一类更一般的数字特征，实际上上面的数字特征都属于特殊的矩。矩作为一类更一般的数字特征，统一描述了随机变量的位置、离散度、形状等信息。

对单个随机变量$X$，它的$k$阶原点矩$\mu'_k$定义为$\E[X^k]$。均值就是$1$阶原点矩。

矩的一个好处是，随机变量的任意多项式函数的均值都可以表示成矩的线性函数。

$k$阶中心矩$\mu_k$定义为$\E[ (X - \E(X])^k )$。方差就是$2$阶中心矩。

$k$阶标准化矩$\tilde{\mu}_k$定义为$E\left( \left( \frac{X - \mu}{\sigma} \right)^k \right)$。$3$阶标准化矩叫做偏度，$4$阶标准化矩叫做峰度。

对随机变量 \(X\) 与 \(Y\)，其 \((m,n)\) 阶或者$m+n$阶混合矩定义为：
\begin{align*}
    &\text{混合原点矩：} \quad \mu_{mn}' = \E[X^m Y^n]\\
    &\text{混合中心矩：} \quad \mu_{mn} = E\left[ (X - \mu_X)^m (Y - \mu_Y)^n \right]
\end{align*}
当 \(m=1, n=1\) 时，$2$阶混合中心矩 \(\mu_{11} = \Cov(X, Y)\) 即协方差。协方差矩阵中的诸元素就是诸$2$阶中心矩。

它真正强大的应用在矩母函数和矩估计处。
